% Zone „kennst du schon" – Kurvenuntersuchung, Kl. 12 Berlin, Gymnasium, GK
\setcounter{aufgabe}{0}
\hypertarget{zone}{}\einheitenkopf{Das kennst du schon}

\begin{aufgabe}{Ich kann besondere Punkte eines Graphen ablesen. Die Abbildung zeigt den Graphen einer Funktion $f$. Lies ab.}

\begin{ksys}[xmin=-1,xmax=5,ymin=-2,ymax=6,ablesen]
\funktionab{\x^3-6*\x^2+9*\x}{f}{-0.19}{4.15}
\end{ksys}

\begin{teile}
\teil Nullstellen von $f$: \leerfeld
\teil Funktionswert $f(2)$: \leerfeld
\teil Koordinaten des Hochpunkts: \leerfeld
\teil Koordinaten des Tiefpunkts: \leerfeld
\teil alle Stellen $x$ mit $f(x) = 4$: \leerfeld
\teil der größte Funktionswert, den $f$ für $0 \le x \le 3$ annimmt: \leerfeld
\end{teile}
\end{aufgabe}

\begin{aufgabe}{Ich kann Gleichungen durch Ausklammern und mit der Lösungsformel lösen. Löse die Gleichung. Schreibe die Umformungen untereinander; die Lösung steht in der letzten Zeile.}
\beispiel{x^2 - 6x + 8 &= 0 \\ x &= 3 \pm \sqrt{9 - 8} \\ x_1 = 4,\ x_2 &= 2}
\begin{gleichungsraster}[2]
\gl{2x - 8 = 0} & \gl{x^2 - 49 = 0} \\
\gl{x^2 - 2x - 8 = 0} & \gl{x^3 - 9x = 0} \\
\gl{x^4 + 3x^3 - 10x^2 = 0} & \gl{(x - 3)\cdot \mathrm{e}^{x} = 0} \\
\gl[3]{x^3 - 2x^2 - 5x + 6 = 0 \text{, eine Lösung ist } x = 1} & \\
\end{gleichungsraster}
\end{aufgabe}

\begin{aufgabe}{Ich finde den Fehler beim Lösen durch Ausklammern. Tom löst die Gleichung $x^3 - 36x = 0$ so:}
\rechnung{x^3 - 36x &= 0 &&\mid :x \\ x^2 - 36 &= 0 &&\mid +36 \\ x^2 &= 36 \\ x = 6 \;&\text{oder}\; x = -6}
\begin{teile}
\teil Finde den Fehler, benenne ihn und korrigiere die Rechnung.
\end{teile}
\begin{gleichungsraster}[3]
\gl{2x^3 - 32x = 0} & \\
\end{gleichungsraster}
\end{aufgabe}

\begin{aufgabe}{Ich kann Ableitungen bilden. Bilde die Ableitung.}
\begin{teile}
\teil $f(x) = x^5$ \feldl{f'(x)}
\teil $f(x) = 3x^2 + 7x$ \feldl{f'(x)}
\teil $f(x) = -0{,}5x^4 + 2x^3 - x + 9$ \feldl{f'(x)}
\teil $f(x) = \frac{1}{3}x^3 - \frac{1}{2}x^2$ \feldl{f'(x)}
\teil $f(x) = x^4 - 5x^3 + x$ \feld{f'(x)} \feld{f''(x)} \feld{f'''(x)}
\teil $f(x) = \mathrm{e}^{-2x}$ \feldl{f'(x)}
\teil $f(x) = (x + 3)\cdot \mathrm{e}^{x}$ \feldl{f'(x)}
\teil $f(x) = x^2\cdot \mathrm{e}^{-x}$ \feldl{f'(x)}
\end{teile}
\end{aufgabe}

\begin{aufgabe}{Ich kann Funktionswerte berechnen und prüfen, ob ein Punkt auf einem Graphen liegt. Berechne den Funktionswert.}
\begin{teilezwei}
\tz $f(x) = x^2 + 1$ \feld{f(3)} & \tz $f(x) = x^3 - x$ \feld{f(2)} \\
\tz $f(x) = 0{,}5x^4 - x^2$ \feld{f(-2)} & \tz $f(x) = -x^2 + 4x$ \feld{f(-3)} \\
\tz $f(x) = x^3 + 2x^2$ \feld{f(-3)} & \tz $f(x) = (x + 1)\cdot\mathrm{e}^{x}$ \feld{f(0)} \\
\end{teilezwei}
Prüfe, ob der Punkt auf dem Graphen von $f$ liegt.
\begin{teile}
\teil $f(x) = x^3 - 5x + 2$, \quad $P(2 \mid 0)$ \quad \janein
\teil $f(x) = x^3 - 5x + 2$, \quad $Q(-1 \mid 4)$ \quad \janein
\teil $f(x) = (x + 1)\cdot\mathrm{e}^{x}$, \quad $R(1 \mid 2\mathrm{e})$ \quad \janein
\end{teile}
\end{aufgabe}

\begin{aufgabe}{Ich kann am Term erkennen, wie ein Graph grob verläuft. Gib an, ob der Graph ganz links und ganz rechts oben oder unten verläuft.}
\begin{teile}
\teil $f(x) = x^2 - 3$ \hfill links \leerfeld rechts \leerfeld
\teil $f(x) = x^3$ \hfill links \leerfeld rechts \leerfeld
\teil $f(x) = -2x^2 + x$ \hfill links \leerfeld rechts \leerfeld
\teil $f(x) = -x^3 + 5x^2$ \hfill links \leerfeld rechts \leerfeld
\teil $f(x) = x^4 - 10x^3$ \hfill links \leerfeld rechts \leerfeld
\end{teile}
Entscheide, ob der Graph symmetrisch zur $y$-Achse (schreibe y), symmetrisch zum Ursprung (schreibe U) oder keines von beiden ist (schreibe k).
\begin{teile}
\teil $f(x) = x^4 - 3x^2$ \leerfeld
\teil $f(x) = x^3 - 7x$ \leerfeld
\teil $f(x) = x^3 + x^2$ \leerfeld
\end{teile}
\end{aufgabe}

\begin{aufgabe}{Ich kann an einem Graphen ablesen, wie steil er ist. Die Abbildung zeigt den Graphen einer Funktion $g$ mit fünf markierten Punkten und der Tangente im Punkt $C$.}

\begin{ksys}[xmin=-4,xmax=4,ymin=-3,ymax=3,ablesen]
\funktionab{(\x^3-12*\x)/8}{g}{-4}{4}
\tangentean*{(\x^3-12*\x)/8}{0}{t}
\punkt{-3}{1.125}{A}
\punkt{-2}{2}{B}
\punkt{0}{0}{C}
\punkt{2}{-2}{D}
\punkt{3}{-1.125}{E}
\end{ksys}

Gib an, ob die Steigung des Graphen im Punkt positiv ist (schreibe $+$), negativ ($-$) oder null ($0$).
\begin{teilezwei}
\tz im Punkt $A$ \leerfeld & \tz im Punkt $C$ \leerfeld \\
\tz im Punkt $B$ \leerfeld & \\
\end{teilezwei}
\begin{teile}
\teil Lies die Steigung der Tangente im Punkt $C$ ab. \feld{m}
\teil Bei welchem der Punkte $B$ und $C$ liegt der Graph höher, bei welchem ist er steiler?
\teil $h(t)$ ist die Höhe eines Ballons in Metern, $t$ Minuten nach dem Start. Es gilt $h'(5) = 2$. Erkläre, was diese Angabe über den Ballon aussagt.
\end{teile}
\end{aufgabe}
